\section{The base map on the original geometric record}

% BEGIN X76R1
\paragraph{X76R1. Level equivariance of the component-normalized reading.}
Fix $T,R,F$ with $\mathrm{Hk3.Core}(R,F)$, a weight and frame tie $w$,
and a pinned base row $O=O_i$. Write $K=T.Q.c.F$,
$M=O.c'.F=\sigma(L)$, and $j=i+1$. Let $C_a$ be the unit
constant of P74I5 for component $a$. For any scalar $u$ the
reading $\Phi_{uC}(f)(h)=u C_{a(h)}\mu(h)
\mathrm{rdU}(f;\mathrm{dict}(h))$ satisfies the original
$\mathrm{EquivAt}$ on lower cusp sections.

Indeed a level lift $k$ of $g\in G$ has trivial infinite part and
preserves the completed reference lattice in both directions. Its
finite action on the marked half is the given reduction
$R.e_G(g)$, hence the endomorphism $O.\mathrm{lev}(g)$ there by the
level pin. The auxiliary level is one. If $\gamma_1,\gamma_2$ are
dictionary movers for $h,hk$, the rational isometry
$\delta=\gamma_2\gamma_1^{-1}$ takes the first marked lattice
to the second with this change of level. This follows by applying
the two lattice and marking identities of the dictionary to
$\gamma_2(hk)=\delta(\gamma_1h)k$. The infinite frame is carried
by $\delta$ since $k_\infty=1$. The resulting isomorphism of
polarized complex tori identifies the two level moves on the
frames, up to the scalar of the theta generator. This scalar
must be kept because the section has exponent $n=\kappa_i.n$.

More explicitly, write $t$ for the ratio of the generator scalars
in compatible period bases, with the second reading
$t^{-1}\vartheta$ in the transported coordinates. The finite
stabilizer acts on the vacuum by $\mu(k)$, as in the stabilizer
calculation of P74I5. The rational motion by $\delta$ changes
variables in the Weil sum. Thus the two descriptions of the
same moved vacuum at an archimedean argument have factors
$m_1=\mu(k)C_{a(h)}\mu(h)$ and $m_2=C_{a(hk)}\mu(hk)$.
The transformed period Gaussian and the $n$ generator readings
give $m_1t^n=m_2$. To cancel the readings, choose one argument
where $\vartheta$ is nonzero (its squared mean norm is positive)
and take it in all $n$ factors. The real period coordinates attain
this argument; the Gaussian is nonzero. Naturality and the
exponent-$n$ rule for an arbitrary section on the same frame now
give
\[
 C_{a(h)}\,\mathrm{rdU}(g f;\mathrm{dict}(h))
       =C_{a(hk)}\,\mathrm{rdU}(f;\mathrm{dict}(hk)),
\]
since the multiplicativity in $w.\mathrm{laws}.\mathrm{muUnit}$ gives
$\mu(hk)=\mu(h)\mu(k)$ and these values are units. Here $gf$ is the
geometric level action. It preserves cusp sections: it preserves
the section and weight laws by change of marking, and on the
finite component charts it transports a cusp reading by the
period isometry just described. Holomorphy and the cusp bound
persist under this change of frame; on unit frames the generator
ratio has unit modulus by the same vacuum comparison. The
argument at every good embedding uses its own polarized
uniformization and normalized generator.

Put $\chi(g)=\prod_{P\in R.P_g}\chi_{j,P}(\det R.e_G(g)_P)$.
The marked determinant satisfies $\chi(g)\mu(k)=1$. To see the
inverse, use a finite-unit determinant root, with residue equal
to the reduced determinant on the marked half and one on the
conjugate half. At a split pair take the determinant at the first
place and one at its mate. At a nonsplit finite place the original
CM extension is unramified there; local Hilbert 90 gives a root
and a power of a base uniformizer makes it a unit. These choices
form a finite idele with trivial infinite component since
$k_\infty=1$. The character is unramified off the marked level
and tame on it. On the marked half its residue character is
$\chi_{j,P}^{-1}$ by the level law. The determinant-root formula
for $\mu$ gives the claimed product. Multiplication of the last
display by $u\mu(h)$ now gives precisely
$\Phi_{uC}(gf)(h)=\chi(g)\Phi_{uC}(f)(hk)$.
% END X76R1

% BEGIN X76R2
\paragraph{X76R2. The base row before the phase is prescribed.}
Continue with the original weak characterization of $R_3,F_3$,
the existence clauses $\mathrm{ExAt}$, the full tie $w$,
$\mathrm{IoExt}$, and the pinned row $O$. These yield the six
core fields and the two tuple character laws by the first
paragraph of T76R1, which uses only the weak characterization
and $\mathrm{ExAt}$. The ordered prime agreement follows from
the half pin of this row by the same contraction proof at $i$.
The base predicate $\mathrm{OneMapRow0}$ holds.

Let $\ell_i=\mathrm{lieFac}>0$, let $c_0$ be the root of unity
from $O.\mathrm{cxi}$, and choose
$u=K_d/(c_0\ell_i)$ as in P74I7 and J74T2. Thus $|u|=1$.
For a definite vector $y$, its coefficient vector $x$ placed
on the prescribed representatives, and any section $f'$ giving
that class sum, their finite theta calculation on every $h$ is
\[
 \mathrm{thF}_{i,\mu_0}(y)(v_0)(h)
   =\ell_i\,u C_{a(h)}\mu(h)\,
             \mathrm{rdU}(f';\mathrm{dict}(h)).
\]
The calculation unfolds the same finite representatives as the
placement of $x$ and is valid before passing to an automorphic
quotient. T73C5 supplies a cusp section for every placed $x$.
With the central group element $1$ and the tuple root of unity
$1\in M$, this gives the full upper row identity: evaluate at
the chosen coset representative, use the pin identifying
$O.\mathrm{thAd}$ there and the identity level action. Its
constants are $uC$ and phase $1$. This construction involves
no division by any lift.

Here is the base change used in the other row fields. Choose a
$K$-basis $(b_j)$ of the number field $M$. For a lower test ring
$S$, evaluate an upper section after extending the framed point
to $M\otimes_K S$ and take its $b_j$-coordinate. The extended
test ring is reduced when $S$ is reduced by separability.
Naturality and the lower weight law commute with these
$S$-linear coordinate projections. The resulting lower sections
$f_j$ give $f'=\sum b_j(f_j)_M$ by expansion in the basis and
the multiplication map on an upper test. For a cusp section
every $f_j$ is cuspidal: at any good embedding of $K$, let $\tau$
run over its extensions to $M$ and let $(b_j^*)$ be trace-dual
to $(b_j)$. All the sums $\sum_k\tau(b_k)\mathrm{rdU}(f_k)$
belong to the required complex cusp subspace on that same
uniformization. Multiply by $\tau(b_j^*)$ and sum over $\tau$.
Trace duality gives $\mathrm{rdU}(f_j)$ in that subspace.
The same restriction of a moved framed point to $K$ and
naturality show that level action commutes with base change.

For a rational base input the placement pin provides $x$ for its
embedded definite vector. Apply the displayed identity to its cusp
section and the preceding decomposition, lifting every $b_j\in M$
back along $\sigma:L\to M$. The unprojected lift thus belongs to
the $\sigma(L)$-span of the functions $\ell_i\Phi_{uC}(f_j)$.
This span is stable under each finite level operator: at the trivial
flag the extra level character is one, and X76R1 expresses right
translation as geometric level action times $\chi(g)^{-1}$,
the image of $(R.FW_i.\chi(g))^{-1}$ by the core group law.
An unavailable level lift gives the zero operator by definition.
The prescribed projector is a finite sum of these operators with
coefficients from $L$. Applying it and the core identity for
$\mathrm{thG}$ gives exactly $O.\mathrm{PhRat}$ for $uC,1$,
including a zero projected lift. The field equality in the pins,
the cusp and action base changes above, the upper row identity
and X76R1 are all the other fields of $O.\mathrm{OneMapOf}(uC,1)$.
Taking these existential witnesses proves the assertion.
% END X76R2

% BEGIN X76R3
\paragraph{X76R3. The guarded all-tuple row on the selected column.}
Take the geometric datum $z$ on the selected column of U76W3.
Assume its weak characterization, $\mathrm{ExAt}$, the original
rank bound, $\mathrm{LUnif}$ and $\mathrm{IoExt}$. Retain the
cover-field conclusion proved in the first paragraph of U76W3.
Then $\mathrm{OneMapRowR}$ with constant phase one holds on every
row of $z$. In particular on this column the base maps in the
last paragraph of U76W3 are outputs.

Indeed $z.\mathrm{geo}$ supplies the tie and pins used in X76R2.
The conclusion of X76R2 is the first conjunct of
$\mathrm{OneMapRowR}$. Apply $z.\mathrm{geo.phase}$ at the
assumed uniformity to this base map to prescribe the constants
$T.c_E$ and the phase one. For the second conjunct suppose the
guard holds. RTG2 gives the adelic character values from the
level and family laws. RTG10 and D74U5 now build the cover
packet on the pinned coefficient field, using the half pin of
this row, the family roots and the retained cover-field conclusion.
RTG12 gives the all-tuple record and RTG4--RTG5 retie it to
the base coefficient witness and level map. T73C5 and P74I7
supply class sums and one common point identity, including after
a character rescaling. F76B2 supplies the base span from the
rank bound and the core and tie just used. The unit-vector
argument of T76R1 gives evaluation injectivity directly from
$w$ and the model law. Consequently RTG16--RTG18 give the
all-tuple $\mathrm{OneMapOf}$ with those prescribed constants
and phase, on a coefficient field equal to $\sigma(L)$. Equality
supplies the inclusion required in the guarded conjunct. This
uses the cover-field route on the selected column. It does not
require a previously assembled $\mathrm{TiesG}$ or its rational
tests. The conditional construction on that column still needs
the uniformity, oscillator extension and geometric datum used
above.
% END X76R3
